% 第14回　キャリア講習 1（記入例・模範解答）
\session{14}{3}{12月16日（水）}{キャリア講習 1}{}{}

\instr{内容や持ち物などの詳細は、決まり次第サイトの各回ページとトップページのお知らせで案内します。}
\instr{\textcolor{ans}{講習の内容が未定のため、記入例は「書き方の例」として示す。}}

\afillin{講習のテーマ・講師}{（例）社会人として働くこと／卒業生の講師（業界・職種を記入）}

\work{メモ}{講習の内容}
\alines{8}{（例）講師の仕事：入社してからの担当業務の変化。最初は営業、現在は企画。\par
・入社前に思っていた仕事と、実際の仕事の違い：「調べる・まとめる・説明する」が毎日ある\par
・学生時代にやっておいてよかったこと：ゼミでの発表、アルバイトでの接客\par
・AIの使い方：資料の下書きには使うが、顧客への説明は必ず自分の言葉でする\par
・就職活動：業界研究は「人に会って聞く」ことで深まった\par
・印象に残った言葉：「わからないことを、わからないと言える人が伸びる」}

\work[60mm]{整理}{学んだこと3つ}
\begin{wstabh}{colspec={Q[c,wd=6mm,m] X[2,l,m] X[3,l,m]},row{2-Z}={ht=16mm}}
  & 学んだこと & 自分のキャリアに引きつけると \\
  1 & \af{仕事の多くは「調べる・まとめる・説明する」の繰り返し} & \af{この授業の前半（問い→文献→確定）と同じ流れ。レポートを丁寧に書くことが練習になる} \\
  2 & \af{AIは下書きまで。説明の責任は自分にある} & \af{第13回の「人間ならではの役割」の定義と重なる。自分の言葉で説明する力を鍛えたい} \\
  3 & \af{業界研究は、人に会って聞くと深まる} & \af{後半の企業との連携で、担当者の一言が一番参考になった経験と同じ} \\
\end{wstabh}

\work{問い}{講師に聞きたいこと・自分で調べたいこと}
\alines{3}{・入社1年目に一番困ったことと、それをどう乗り越えたか。\par
・文系の学生が、今のうちに身につけておくとよい力は何か。\par
・（自分で調べる）講師の業界の主な企業と、仕事の種類。}

\begin{nexttime}
  \item 第15回（12月23日）に向けて、配布資料があれば持参する
\end{nexttime}

\begin{point}
  \item 講習の内容が決まり次第、記入例を差し替える。ここでは、メモの取り方と「自分に引きつける」書き方の目安を示した。
  \item 「学んだこと」と「自分のキャリアに引きつけると」を分けて書かせる。後者に、この授業の経験（問いの立て方、AIとの協働、企業との連携）を結びつけられていれば高く評価する。
  \item 講師への質問は、講習中に1つは書かせ、質疑の時間に使わせる。
\end{point}
